\begin{papercard}{artstein2008}
\cf{Stages} \stagetag{I} (also \stagetag{C} for unit and category design, and \stagetag{L} for the caveat that agreement does not equal machine-learnability). A 41-page survey from computational linguistics (CL) of how annotator agreement coefficients work (unified notation, mathematics, assumptions), of the ``bias'' and ``prevalence'' debates, and of ten years of use on four annotation tasks. It argues for weighted, $\alpha$-like coefficients but concedes that these are harder to interpret. Only part of the material fits in the article; an extended online version exists.

\cf{Reliability, agreement, validity} Hand-coded data must be shown reliable: coders agree on the categories to an extent set by the purpose of the study. Agreement indicates that coders have internalised a similar reading of the guidelines. Reliability is a prerequisite for validity but cannot ensure it, since two observers may share the same prejudice \pc{artstein2008}{3}{However, it is important to keep in mind that achieving good agreement cannot ensure validity}. Notation: items $I$, categories $K$, coders $C$; observed/expected agreement $A_o$, $A_e$; observed/expected disagreement $D_o$, $D_e$.

\cf{Percent agreement fails twice} Observed agreement is the mean of per-item agreement ($A_o=.7$ in the dialogue-act example, 100 utterances, 2 coders). It cannot be compared across studies because chance agreement depends on (i) the \emph{number of categories} (a binary scheme gives 1/2 by chance, a three-category scheme 1/3) and (ii) the \emph{distribution} of items. With 95\,\% statements, chance already yields 90.5\,\% agreement, so ``a seemingly high observed agreement of 90\,\% is actually worse than expected by chance'' \pc{artstein2008}{5}{Under such circumstances, a seemingly high observed agreement of 90\% is actually worse than expected by chance}.

\cf{Common formula and three chance models} $S,\pi,\kappa=(A_o-A_e)/(1-A_e)$, with value 0 at chance and 1 at perfect agreement. All three assume the coders are independent; they differ only in $P(k\mid c)$: $S$ assumes a uniform distribution over categories, $\pi$ one common distribution estimated from all judgments, $\kappa$ a separate distribution per coder \pc{artstein2008}{6}{S: This coefficient is based on the assumption that if coders were operating by chance alone, we would get a uniform distribution}. On Table 1 (observed .7): $S=.4$, $\pi\approx.341$, $\kappa\approx.348$. $S$ can be inflated by adding categories the coders never use, so it would reward over-fine tagsets \pc{artstein2008}{7}{The value of the coefficient can be artificially increased simply by adding spurious categories which the coders would never use}. Always $\pi\le\kappa$, with equality when the coders' distributions are identical.

\cf{More than two coders} Averaging pairwise values is common, but ``a better practice is to use generalized versions of the coefficients'': Fleiss' multi-$\pi$ (misleadingly called ``kappa'' in the literature; Siegel \& Castellan's $K$ is the same coefficient) and Davies \& Fleiss' multi-$\kappa$ \pc{artstein2008}{8}{However, a better practice is to use generalized versions of the coefficients}. Agreement per item is the proportion of agreeing judgment pairs; an agreement table (items $\times$ categories) replaces the contingency table, which loses the information about \emph{which} coder gave which label.

\cf{Krippendorff's $\alpha$ and weighted $\kappa$} Neither $\pi$ nor $\kappa$ distinguishes between kinds of disagreement. $\alpha$ is derived from analysis of variance ($1-s^2_{\text{within}}/s^2_{\text{total}}$) and then generalised to $\alpha=1-D_o/D_e$ with any distance function $d$ that is zero for identical labels; it handles multiple coders, missing values and any measurement level \pc{artstein2008}{10}{A serious limitation of both}. For nominal distance ($d=0/1$) $D_o$ equals one minus the multi-$\pi$ observed agreement and $D_e$ differs only by the factor $(ic-1)/ic$, so $\alpha\approx\pi$ and the two coincide as items or coders grow. The flexibility has a price: different distance metrics give different $\alpha$ for the same data and ``it is not clear what the natural distance metric is'' \pc{artstein2008}{13}{the flexibility also creates new pitfalls, especially in cases where it is not clear what the natural distance metric is}. Cohen's $\kappa_w$ uses weights per category pair and individual coder distributions. Integrated example (Table 4, three categories, $A_o=.88$): $S=.82$, $\pi=.7995$, $\kappa=.8013$; with weights 0/0.5/1, $D_o=.09$ and $\alpha\approx.8156$, $\kappa_w=.8163$.

\cf{Bias and prevalence} \emph{Bias}: Di Eugenio \& Glass said $\kappa$ is paradoxical because it rises as coders become less alike, and recommended reporting both $\kappa$ and $\pi$. The authors answer that this misunderstands the difference: $\pi$/$\alpha$ estimate the chance of agreement among \emph{arbitrary} coders, $\kappa$ among \emph{these} coders, so the choice depends on the intended interpretation, not on how divergent the marginals look \pc{artstein2008}{17}{Therefore, the choice of coefficient should not depend on the magnitude of the divergence between the coders}. $\pi$ measures reproducibility; $\kappa$ ``rewards'' biased annotators as a question of trustworthiness (validity) \pc{artstein2008}{17}{it is not a matter of reproducibility (reliability) but rather of trustworthiness (validity)}. Overall annotator bias $B=A^\pi_e-A^\kappa_e=\frac{1}{c-1}\sum_k\sigma^2_{\hat P(k\mid c)}$ and $\kappa-\pi=B(1-A_o)/[(1-A^\kappa_e)(1-A^\pi_e)]$, so the difference shrinks with more coders and higher agreement (18 coders: $\alpha$ and a $\kappa$-style $\alpha$ differed by less than 0.001). Advice: choose the coefficient by task rather than by observed marginals; for reliability use $\pi$/$\alpha$; test with more than two coders \pc{artstein2008}{18}{the appropriate coefficient should be chosen based on the task}. \emph{Prevalence}: with 95\,\% in one category random coding already gives 90.25\,\% joint assignment and at least 85.7\,\% accuracy, yet coefficients stay low, which is ``justified'' because reliability means the ability to discriminate; ``the test for reliability in such cases is the ability to agree on the rare categories'' \pc{artstein2008}{19}{The test for reliability in such cases is the ability to agree on the rare categories}. With margins equal and rare-category agreement $\delta$, half-disagreement $\epsilon$: $\pi=\kappa\approx\delta/(\delta+\epsilon)$.

\cf{Krippendorff's methodology as applied to CL} Reproducibility data need (i) an exhaustive coding scheme with step-by-step instructions, (ii) specified coder-selection criteria, (iii) independent coders \pc{artstein2008}{20}{It must ensure that the coders that generate the data used to measure reproducibility work independently of each other}. CL violates these by extending instructions as data are generated, by using long-term collaborators as coders, and by letting coders discuss and vote \pc{artstein2008}{21}{violated by using experts as coders, particularly long-term collaborators}. Three or more coders; the sample must include every category often enough for at least five chance agreements; examples of all values and domains \pc{artstein2008}{21}{each category should occur in the sample often enough to yield at least five chance agreements}. \emph{Significance}: testing against chance is uninteresting (Cohen); what matters is a confidence interval, obtained for $\kappa$ with normal theory and for $\alpha$ by bootstrapping (Hayes \& Krippendorff 2007) \pc{artstein2008}{22}{confidence intervals must be obtained by bootstrapping}.

\cf{Interpretation} Landis \& Koch bands are used in medicine; CL adopted Krippendorff's $.8$ good and $.67$--$.8$ tentative (``highly tentative and cautious'' originally), though Krippendorff later called $\alpha=.8$ the absolute minimum; Neuendorf's $\ge.9$ and $\ge.8$. Deciding what counts as adequate is ``little more than a black art'' \pc{artstein2008}{22}{deciding what counts as an adequate level of agreement for a specific purpose is still little more than a black art}. Reidsma \& Carletta showed that agreement coefficients are poor predictors of machine-learning success; reliability and machine-learnability are different questions \pc{artstein2008}{23}{agreement coefficients should not be used as indicators of the suitability of annotated data for machine learning}.

\cf{Annotation tasks} \emph{Dialogue acts}: flat schemes reach $K\approx.8$--$.9$ (MapTask 13 tags .83; Switchboard-DAMSL 42 tags .80; CSTAR 20 tags .90) but assume mutually exclusive categories; multi-dimensional DAMSL gave $K$ from .15 to .76 (Core \& Allen) and .54 to .93 (Di Eugenio et al.); hierarchical DIT++ with $\kappa_w$ gives higher values that are not directly comparable. \emph{Segmentation and boundaries}: $K$ between .36 and .97, mostly low because the coders agree on the bulk but not the exact boundary; with the boundary class rare, expected agreement is high, so two coders disagreeing on one boundary of 50 utterances get $K=.65$ whereas a finer scheme with three times as many disagreements gets $.75$ \pc{artstein2008}{27}{the value of K is fairly low}. \emph{Unitizing}: agreement on the identification of units is rarely checked; Krippendorff's $\alpha_U$ squares the length of non-overlapping stretches and assumes non-overlapping units \pc{artstein2008}{28}{The one proposal for measuring agreement on markable identification}. \emph{Anaphora}: labels are sets (coreference chains); Passonneau's set distance (0, 1/3, 2/3, 1), Jaccard, Dice and MASI ($d_P\times d_J$) give $\alpha=.642$--$.691$, while $K$ with chains as labels collapses to .48 because nobody marks all mentions \pc{artstein2008}{31}{the value of K was drastically reduced, because hardly any coder identified all the mentions of discourse entities}. \emph{Discourse deixis}: the same data give $\alpha$ from .122 to .998 depending on how agreement is measured, so rules of thumb no longer apply \pc{artstein2008}{32}{we cannot anymore interpret the value of}. \emph{Word senses}: naive coders on fine-grained dictionary senses reach $K\approx.36$--.67 ($\kappa_w$ .41--.46 for V\'eronis); grouped senses raise percent agreement from 71\,\% to 82\,\%, but post hoc merging is not equivalent to a study with fewer categories to begin with \pc{artstein2008}{34}{post hoc merging of categories is not equivalent to running a study with fewer categories to begin with}. Melamed \& Resnik's hierarchical agreement is recast as the distance $d=1-\sum_l\min(P(l\mid A),P(l\mid B))$ for $\alpha$.

\cf{Recommendations} (1) Perform and report rigorous reliability testing, not percent agreement alone; detailed calculations show where disagreement lies \pc{artstein2008}{36}{Our first recommendation is that annotation efforts should perform and report rigorous reliability testing}. Expert coders are acceptable in CL (usefulness versus replicability) but then results are not replicable across sites. (2) $K$ or $\alpha$ for reliability; annotator bias matters little with many coders, so \emph{increase the number of annotators}; $\alpha$ is indispensable for hierarchical or set-valued labels, but the weights affect the value. (3) Interpretation: 0.8 is a good threshold if one is needed, 0.7 can suffice for some discourse corpora; a threshold should not be a prerequisite for publication; instead report number of coders, independence, reliance on the manual, significance, and a confusion matrix or agreement table so that readers can see whether overall figures hide disagreement on rare categories \pc{artstein2008}{37}{provide a confusion matrix or agreement table so that readers can find out whether overall figures of agreement hide disagreements on less common categories}. Corpus quality should also be judged by independent evaluation of the downstream task.

\cf{Limits} A survey; the author-preferred position on $\pi$/$\alpha$ over $\kappa$ is argued, not neutral; most task examples (dialogue acts, anaphora, senses) are linguistic and do not directly map onto privacy features; thresholds are experience-based; the arithmetic for weighted $\alpha$ depends on distance metrics for which no standard exists.

\cf{Use for this thesis} (1) A uniform notation and the single formula for $S$, $\pi$, $\kappa$ and $\alpha$ so that the methodology can explain, in one paragraph, why the chosen coefficient is used and what the others would show. (2) The rule that the choice follows the intended interpretation (reproducibility of the codebook by other coders $\Rightarrow$ pooled-distribution coefficients such as $\alpha$), not the observed marginals, and that two coders are the minimum, more coders reduce the $\pi$/$\kappa$ gap. (3) Ordinal or hierarchical codes (e.g.\ severity levels, a dark-pattern taxonomy) need an explicitly justified distance metric because the value of $\alpha$ moves with it; record the metric before coding. (4) For rare codes, report the agreement table and agreement on the rare category itself, not only an overall figure. (5) Decide whether unit identification (which screen, setting or policy sentence is coded) is part of the reliability claim, and if so measure it (unitizing) separately. (6) For stage L the paper's reminder that agreement with a human gold standard says nothing about validity, and that low agreement with systematic disagreement can still be learnable. The 0.8/0.7 thresholds here are CL conventions and need to be set against Krippendorff (2004, 2019) and Hallgren (2012).
\end{papercard}
